\documentclass[11pt]{article}
\usepackage[margin=1.15in]{geometry}
\usepackage{amsmath,amssymb,amsthm,mathtools}
\usepackage{xcolor}
\usepackage[colorlinks=true,linkcolor=blue!60!black,citecolor=blue!60!black,urlcolor=blue!60!black]{hyperref}

\newtheorem{theorem}{Theorem}[section]
\newtheorem{proposition}[theorem]{Proposition}
\newtheorem{lemma}[theorem]{Lemma}
\newtheorem{corollary}[theorem]{Corollary}
\theoremstyle{definition}
\newtheorem{definition}[theorem]{Definition}
\theoremstyle{remark}
\newtheorem{remark}[theorem]{Remark}

\newcommand{\C}{\mathbb{C}}
\newcommand{\R}{\mathbb{R}}
\newcommand{\Z}{\mathbb{Z}}
\newcommand{\w}{\omega}
\newcommand{\Sp}{\mathrm{Sp}}
\newcommand{\dg}{\dagger}
\newcommand{\WF}{\operatorname{WF}}
\newcommand{\tr}{\operatorname{tr}}

\title{The Universal Vacuum of the Fourth-Order Scalar Field:\\
Metric Orbits, Fock Sectors, and the Krein Boundary}
\author{}
\date{}

\begin{document}
\maketitle

\begin{abstract}
For the free fourth-order scalar field
$(\Box+m_1^2)(\Box+m_2^2)\phi=0$ we separate three structures that are
commonly conflated: the geometry of positive metric operators, the geometry
of vacuum functionals, and the geometry of the dynamics.  Modewise, the
Bender--Mannheim generator is a pure beam splitter in canonical
normal-form coordinates --- it carries divergent Cartan distortion while
creating no quanta there --- and no symplectic frame tames its singular
values: particle content is a vacuum-ray functional, not a Cartan-norm
functional.  The entire family of admissible positive metrics fixes the
physical vacuum ray, so the vacuum functional is metric-independent; it
carries
exactly $1/3$ of a quantum per ultraviolet mode pair relative to the
two-field Fock vacuum, and its GNS representation is disjoint from the
latter in every spatial dimension --- an obstruction no choice of metric
can cure.  All nondegenerate selected vacua lie in one universal
ultraviolet-dressed Fock sector for $d\le3$, anchored by the $\Box^2$
vacuum, with a terminating sector hierarchy labeled by $m_1^2+m_2^2$ for
$4\le d<8$ and by the full mass pair for $d\ge8$.  The selected vacuum's
Wightman function is exactly the spectral two-point function of the
fourth-order operator for every mass splitting, and its equal-mass limit
coincides with the Krein-space vacuum of Bateman and Turok: metric
selection and the spectral condition determine the same quasifree vacuum
functional (after matching the action sign and infrared extension), the
two quantizations differing only in the completion --- positive
pseudo-Hermitian for split masses, Krein/Jordan at resonance.  Finally the
selected functional is Hadamard in the natural fourth-order sense: its
wavefront set is the standard positive-frequency Hadamard set, with a
logarithmic short-distance singularity of universal coefficient
$1/(16\pi^2)$ for $\log\rho^2$, and its restriction to the local
partial-fraction fields is $(\pm)$Klein--Gordon--Hadamard.  All identities
are machine-verified.
\end{abstract}

\section{Introduction}\label{sec:intro}

The Pais--Uhlenbeck (PU) oscillator and its field-theoretic parent, the
free fourth-order scalar
\begin{equation}\label{eq:eom}
(\Box+m_1^2)(\Box+m_2^2)\,\phi=0,\qquad m_1>m_2\ge0,
\end{equation}
admit two apparently different resolutions of the Ostrogradsky ghost
problem.  In the PT-symmetric approach of Bender and Mannheim
\cite{BM2008PRL,BM2008PRD}, each mode carries a positive metric operator
$\eta=e^{-Q}$ that renders the dynamics unitary on a Hilbert space; in the
Krein-space approach, recently developed for the pure fourth-order theory
$\Box^2\phi=0$ by Bateman and Turok \cite{BT2026}, the state space is
indefinite and probabilities are controlled by a hidden ghost-parity
symmetry.  Two companion papers \cite{companion1,companion2} analyzed the
modewise PT construction as finite-dimensional symplectic geometry: the
positive metric is unique up to a stabilizer orbit, is selected
variationally by minimum Cartan distortion, and the modewise theory is
exactly solvable in all the respects used below.

This paper carries that analysis to the field theory, where three
different geometries are in play and must not be conflated:
\begin{equation}\label{eq:threegeom}
\text{metric geometry}\ \neq\ \text{vacuum geometry}\ \neq\
\text{dynamical geometry}.
\end{equation}
The metric geometry is the symmetric-space geometry of positive operators,
$G/K$ for $G=\Sp(2n,\C)$, where the Cartan projection $\mu$ measures
nonunitarity \cite{companion2}.  The vacuum geometry is the geometry of
Gaussian vacuum rays, which is governed by a different quotient: the
stabilizer $P_\Omega=\{g:\hat g\,\Omega\in\C^\times\Omega\}$ of the
vacuum line is a parabolic-type subgroup, and vacuum rays live in
$G/P_\Omega$.  There is no map $G/K\to G/P_\Omega$ (the compact
$K\cong\Sp(n)$ does not preserve the vacuum ray; only $K\cap
P_\Omega\cong U(n)$ does); the correct geometry is the correspondence
$G/K\leftarrow G\rightarrow G/P_\Omega$ of two inequivalent quotient
maps.  The dynamical geometry is the spectral structure of the
Hamiltonian, which degenerates (Jordan blocks) exactly at the resonant
point $m_1=m_2$ \cite{companion1}.  Our results, in brief:

\begin{enumerate}
\item \emph{Cartan--parabolic decoupling}
(Section~\ref{sec:decoupling}).  In
canonical normal-form coordinates the Bender--Mannheim generator is a pure
beam splitter, $Q'=r\,(a_1a_2^{\dg}+a_1^{\dg}a_2)$: it lies in the
complexified number-preserving (Levi) algebra of $P_\Omega$, fixes the
canonical vacuum exactly, and yet has Cartan projection $(r,r)$ with
$r=\log\frac{\w_1+\w_2}{\w_1-\w_2}\to\infty$ in the ultraviolet.
Moreover no symplectic frame tames the singular values.  Hence there is
no inequality of the form $\text{vacuum excitation}\ge f(\|\mu\|)$:
Cartan distortion does not control particle content.

\item \emph{Orbit constancy and a universal obstruction}
(Section~\ref{sec:obstruction}).  The physical annihilator covectors are
eigenvectors of the stabilizer generators, so the entire admissible
metric family --- the Gaussian orbit and all $W=f(N_1,N_2)>0$ ---
fixes the physical vacuum ray.  The vacuum functional is therefore
metric-independent.  Its occupation relative to the two-field Fock vacuum
is exactly
\begin{equation}\label{eq:N}
\langle N\rangle(k)
=\frac{\w_2(\w_1^2+\w_2^2)}{2\w_1(\w_1^2+\w_1\w_2+\w_2^2)}
\ \xrightarrow[k\to\infty]{}\ \frac13 ,
\end{equation}
a universal constant; the total relative quanta scale as
$\Theta(V\Lambda^d)$ and the representations are disjoint in every
spatial dimension.  No choice of metric can cure this.

\item \emph{One universal dressed sector; terminating hierarchy}
(Section~\ref{sec:sectors}).  Selected vacua of different mass pairs
satisfy, in common field variables,
\begin{equation}\label{eq:sector}
1-f(k)=\frac{(\Sigma-\bar\Sigma)^2}{12\,k^4}+O(k^{-6}),
\qquad \Sigma=m_1^2+m_2^2 ,
\end{equation}
with no $k^{-6}$ term and next obstruction
$\tfrac{29}{576}\,(\delta^2-\bar\delta^2)^2k^{-8}$,
$\delta=\tfrac12(m_1^2-m_2^2)$.  Hence: one universal sector for
$d\le3$, anchored by the $\Box^2$ vacuum; sectors labeled by $\Sigma$
for $4\le d<8$; by the full mass pair for $d\ge8$; and the
classification terminates.

\item \emph{The spectral bridge} (Section~\ref{sec:bridge}).  The
selected vacuum's Wightman function is exactly the spectral two-point
function of \eqref{eq:eom},
\begin{equation}\label{eq:Wspectral}
\widetilde W^+_{12}(p)
=\theta(p^0)\,
\frac{\delta(p^2-m_1^2)-\delta(p^2-m_2^2)}{m_1^2-m_2^2},
\qquad
\widetilde W^+_{mm}(p)=-\theta(p^0)\,\delta'(p^2-m^2),
\end{equation}
for all $\Delta=m_1^2-m_2^2\ge0$; the confluent case is the Krein vacuum
of \cite{BT2026}.  Metric selection and the spectral condition determine
the same quasifree vacuum functional; the quantizations differ only in
the completion --- positive $\eta$-inner product where it exists
($\Delta>0$), Krein space with ghost parity where it must ($\Delta=0$).
The state is analytic across the resonant boundary while the metric and
diagonalizability fail there.

\item \emph{Fourth-order Hadamard property} (Section~\ref{sec:hadamard}).
$W^+_{12}$ is a bisolution with the correct fourth-order commutator, its
wavefront set is the standard positive-frequency Hadamard set
$\mathcal C^+$, and its short-distance singularity is logarithmic with
universal coefficient: the $1/\rho^2$ singularity of the Klein--Gordon
factors is mass-independent and cancels, leaving
\begin{equation}
W^+_{12}(x)=\frac{1}{8\pi^2}\log\rho+\text{(smooth)},
\end{equation}
i.e.\ coefficient $1/(16\pi^2)$ for $\log\rho^2$.  Its restrictions to
the local partial-fraction fields are exactly $(\pm)$Klein--Gordon
Hadamard functions.  Wick powers and perturbation theory can be based on
the fourth-order parametrix
$H_P=(H_{m_1}-H_{m_2})/(m_1^2-m_2^2)$.
\end{enumerate}

Together these give the corrected physical picture: \emph{the universal
selected vacuum is a fourth-order Hadamard quasifree functional; its
apparently infinite particle content relative to the split two-field
vacuum reflects inequivalent complex structures, not pathological local
ultraviolet singularities} --- and the Krein formulation is not the
remnant of a divergent positive theory but a different completion of the
same, analytically continued vacuum.

\paragraph{Verification.}
Every identity below --- the beam-splitter identity, the occupation
\eqref{eq:N}, the eigenvector/constancy lemma, the sector expansions with
their exact coefficients, the Wightman functions and their confluent
limits, the commutator normalizations, and the logarithmic coefficient ---
is machine-verified symbolically, with independent numerical checks; the
scripts accompany the repository of \cite{companion1,companion2}.

\paragraph{Conventions.}
Modewise notation follows \cite{companion1,companion2}: for spatial
momentum $k$, the mode is a PU pair with
$\w_i(k)=\sqrt{k^2+m_i^2}$, canonical variables $V=(x,y,p,q)$ related to
the field variable by $z=iy$ (the PT rotation), normalized coordinates
with $d_xd_y=\w_1\w_2$, rapidity $r(k)=\log\frac{\w_1+\w_2}{\w_1-\w_2}
=\log\frac{(\w_1+\w_2)^2}{\Delta}$, selected diagonalizer
$S_+=e^{(r/2)B}$ and metric $\eta=e^{-Q}$, $Q=\alpha pq+\beta xy$.  We
write $\Delta=m_1^2-m_2^2$, $\Sigma=m_1^2+m_2^2$, $\Pi=m_1^2m_2^2$.

\section{Metric geometry versus vacuum geometry}\label{sec:decoupling}

\subsection{The beam-splitter identity}

In the canonically normalized coordinates the mode annihilation operators
of the auxiliary canonical vacuum are $a'_j$ with
$x'=(a'_1+a'^{\dg}_1)/\sqrt2$, etc.  A one-line computation (different
modes commute, so no ordering terms arise) gives the exact operator
identity
\begin{equation}\label{eq:beamsplitter}
Q'=r\,(x'y'+p'q')=r\,\big(a'_1a'^{\dg}_2+a'^{\dg}_1a'_2\big).
\end{equation}
Thus $Q'$ is a pure beam splitter: $[Q',N_{\mathrm{tot}}]=0$ and
$Q'\,\Omega_{\mathrm{can}}=0$, so
\begin{equation}
e^{-Q'/2}\,\Omega_{\mathrm{can}}=\Omega_{\mathrm{can}}
\qquad\text{exactly},
\end{equation}
even though the Cartan projection of the implemented symplectic map is
$\mu(S_+)=(r,r)$ with $r(k)\sim2\log k\to\infty$.  Equivalently: $Q'$
lies in the complexified number-preserving algebra, the Levi component of
the parabolic-type stabilizer $P_\Omega$ of the vacuum line, on which the
Cartan norm is unbounded.

\subsection{No frame tames the map; the occupation is a state functional}

One might hope that a frequency-adapted symplectic frame $T_\w$ (the one
whose canonical vacuum is the physical two-field vacuum) turns the
divergent Cartan data into bounded physical squeezing.  It cannot:

\begin{lemma}[No frame tames the map]\label{lem:noframe}
For every invertible frame $T$,
\begin{equation}
\big\|\mu\big(T^{-1}S_+T\big)\big\|_\infty
\ \ge\ \log\varrho\big((T^{-1}S_+T)^2\big)=\log\varrho(S_+^2)=r ,
\end{equation}
where $\varrho$ is the spectral radius: the largest singular value of any
matrix dominates its spectral radius, and similarity preserves the
spectrum of the positive $S_+$, which is $\{e^{\pm r/2}\}$.  Hence
$\|\mu(T_\w^{-1}S_+T_\w)\|\to\infty$ in the ultraviolet for every choice
of frames $T_\w$.  (Numerically the physical-frame Cartan data are in
fact close to $(r,r)$; the lemma is what the theorems below use.)
\end{lemma}

What is bounded is the \emph{state}: the selected vacuum, computed exactly
in \cite{companion2} as the Gaussian ground state of the PT mode
Hamiltonian, has occupation \eqref{eq:N} relative to the two-field vacuum,
saturating at $1/3$.  The two statements are compatible because particle
content depends on the transformation only through the vacuum ray it
produces --- a $G/P_\Omega$ datum --- while $\mu$ is a $G/K$ datum, and
the two are inequivalent quotients of $G$.

\begin{theorem}[Cartan--parabolic decoupling]\label{thm:decoupling}
Let $G=\Sp(2n,\C)$ act metaplectically on Fock space, and let
$q_K:G\to G/K$ and $q_\Omega:G\to G/P_\Omega$,
$q_\Omega(g)=[\hat g\,\Omega]$, be the two quotient maps.  Then:
(i) positive-metric geometry is a function of $q_K$ (the Cartan
projection); (ii) Gaussian vacuum rays are a function of $q_\Omega$;
(iii) there is no map $G/K\to G/P_\Omega$ intertwining them, since
$K\not\subseteq P_\Omega$ ($K\cap P_\Omega\cong U(n)\subsetneq
K\cong\Sp(n)$); (iv) the Cartan norm is unbounded on $P_\Omega$ (already
on its Levi component), so fibres of $q_\Omega$ contain elements of
arbitrarily large Cartan norm; consequently (v) no bound of the form
$\langle N\rangle\ge f(\|\mu\|)$ or $\le f(\|\mu\|)$ holds in general.
The PU metric orbit realizes the extreme case:
$q_\Omega(S_+H)=\{[\Omega_{\mathrm{PT}}]\}$ is a single vacuum ray
(Theorem~\ref{thm:constancy}), while $q_K(S_+H)$ is unbounded, with
$\mu(S_+)=(r,r)$ divergent in the ultraviolet yet identically vanishing
canonical-frame excitation, and bounded physical excitation \eqref{eq:N}
despite Lemma~\ref{lem:noframe}.
\end{theorem}

The proof of (iv) is the identity \eqref{eq:beamsplitter} together with
$\mu(e^{(r/2)B})=(r,r)$; (v) follows from the exhibited extremes.  This
corrects a natural but false expectation --- used, e.g., to estimate
Bogoliubov costs by $\sinh^2$ of Cartan coordinates --- and it is the
conceptual reason the theorems below must be formulated at the level of
vacuum functionals.

\section{Orbit constancy and the universal Fock obstruction}
\label{sec:obstruction}

\begin{lemma}[Eigenvector lemma]\label{lem:eigen}
Let $\ell_j$ be the physical annihilator covectors (the two-field vacuum
modes) and $X_j$ the stabilizer generators of \cite{companion2}.  Then
\begin{equation}
\ell_j^{\top}X_j=-i\,\ell_j^{\top},
\qquad\text{hence}\qquad
\ell_j^{\top}C(\theta_1,\theta_2)^{-1}=e^{i\theta_j}\,\ell_j^{\top}
\quad(\theta_j\in\C).
\end{equation}
\end{lemma}

\begin{theorem}[Metric-independence of the vacuum]\label{thm:constancy}
The selected physical vacuum ray, and hence the associated quasifree
functional, is the same for every admissible positive metric: (i) on the Gaussian orbit $\eta_C$,
$C\in SO(2,\C)^2$, the annihilation conditions defining the vacuum are
unchanged by Lemma~\ref{lem:eigen}; (ii) the metaplectic implementers of
the stabilizer are of the form
$\exp[\theta_1(N_1+\tfrac12)+\theta_2(N_2+\tfrac12)]$ and act on the
vacuum ray by scalars; (iii) for the full metric family
$\eta'=\rho^{\dg}W\rho$ with $W=f(N_1,N_2)>0$ \cite{companion1}, $W^{\pm1/2}$
also fixes the vacuum ray.  Hence $[\Omega_{\eta}]=[\Omega_{\eta_+}]$ for
every admissible $\eta$.
\end{theorem}

The theorem asserts constancy of the vacuum \emph{ray}; metric changes do
alter the physical inner product and may alter representatives of
observables, but any quantity determined by the fixed vector relative to
the fixed two-field reference --- in particular the occupation
\eqref{eq:N} and the fidelity below --- is metric-independent.

\begin{theorem}[Universal Fock obstruction]\label{thm:nogo}
For every admissible positive metric, in a volume $V$ with cutoff
$\Lambda$: the relative quanta obey
$N(\Lambda)\sim C_d\,V\Lambda^d$ with $C_d>0$; the vacuum fidelity obeys
$|\langle\Omega_{2\text{-field}}|\Omega_{\mathrm{PU}}\rangle|^2
\le e^{-cV\Lambda^d}$ (per-mode fidelity $\to\sqrt3/2$
\cite{companion2}); and in the thermodynamic/continuum limit the GNS
representations are disjoint, in every spatial dimension $d\ge1$.
Metric ambiguity cannot cure the inequivalence.
\end{theorem}

\begin{proof}
Combine Theorem~\ref{thm:constancy} with the exact modewise data
\eqref{eq:N} and the von Neumann infinite-tensor-product criterion, as in
\cite{companion2}.
\end{proof}

\section{The universal dressed sector}\label{sec:sectors}

Comparisons between selected vacua of different mass pairs are made in
common physical variables ($\gamma=1$ PU variables per mode; the
normalization map contributes only $O(k^{-2})$).  The exact modewise
Gaussian of \cite{companion2} gives, for pairs $(m_1,m_2)$ and
$(\bar m_1,\bar m_2)$, the fidelity expansion \eqref{eq:sector}; at
matched $\Sigma$ the expansion continues
\begin{equation}\label{eq:sector8}
1-f(k)\Big|_{\Sigma=\bar\Sigma}
=\frac{29}{576}\,\frac{(\delta^2-\bar\delta^2)^2}{k^{8}}+O(k^{-10}),
\qquad \delta^2-\bar\delta^2=\bar\Pi-\Pi ,
\end{equation}
with \emph{no} $k^{-6}$ term.

\begin{theorem}[Sector classification]\label{thm:sectors}
Let $d$ be the spatial dimension.  Relative Hilbert--Schmidt convergence
$\int^\Lambda k^{d-1}(1-f(k))\,dk<\infty$ holds iff the first nonvanishing
invariant difference is compatible with $d$:
\begin{center}
\begin{tabular}{c|c}
spatial dimension & UV sector invariant of selected vacua\\
\hline
$d\le3$ & none (one universal sector)\\
$4\le d<8$ & $\Sigma=m_1^2+m_2^2$\\
$d\ge8$ & $(\Sigma,\Pi)$, i.e.\ the full mass pair
\end{tabular}
\end{center}
The hierarchy terminates: two symmetric invariants determine the pair.
At the critical dimensions the obstruction is logarithmic rather than
power-law: at $d=4$ (respectively $d=8$) the divergent sum for
$\Sigma\neq\bar\Sigma$ (respectively $\Pi\neq\bar\Pi$ at matched
$\Sigma$) grows as $\log\Lambda$, so the critical dimensions belong to
the labeled regimes as indicated.
In $d\le3$ the universal sector is anchored by the $\Box^2$ selected
vacuum ($m_1=m_2=0$; modewise Gaussian $A=(2k,k^2;k^2,2k^3)$), relative to
which $1-f=\Sigma^2/(12k^4)+\dots$; all massive selected vacua are
quasi-equivalent to it and to each other, while remaining collectively
disjoint from the naive two-field sector by Theorem~\ref{thm:nogo}.
\end{theorem}

At the resonant boundary three kinds of regularity must be
distinguished.
\emph{Distributional regularity}: the selected two-point functional is
analytic in $\Delta$ across $\Delta=0$ (modewise, $1-f\sim\varepsilon^4$
for $m_{1,2}=m\pm\varepsilon$; the Wightman functions of
Section~\ref{sec:bridge} converge smoothly), and the occupation limits
commute ($\langle N\rangle\equiv\tfrac13$ at $\Delta=0$ for every $k$).
\emph{Representation regularity}: for $d\le3$ the resonant vacuum lies in
the same dressed Fock sector as every nondegenerate selected vacuum
(Theorem~\ref{thm:sectors}).
\emph{Dynamical and metric singularity}: as $\Delta\to0$ the metric
operator degenerates ($\alpha,\beta\to\infty$), the similarity
transformation diverges ($S_+\sim\Delta^{-1/2}$), and the dynamics becomes
nondiagonalizable (Jordan blocks; no positive invariant form
\cite{companion1}).  In summary,
\begin{equation}
\boxed{\ \text{resonance is not a vacuum singularity; it is a metric and
dynamical singularity.}\ }
\end{equation}

\section{The spectral bridge to the Krein quantization}\label{sec:bridge}

The physical Wightman function of the selected vacuum follows from the
exactly solvable mode: with
$\widetilde V=S_+V$ evolving under the diagonalized Hamiltonian,
$W(t)=S_+e^{tA_0}\Gamma S_+^{\top}$ pulled back to the field variable
$z=iy$ gives, exactly,
\begin{equation}\label{eq:Wmode}
W_{zz}(t)
=\frac{1}{\Delta}
\left[\frac{e^{-i\w_1t}}{2\w_1}-\frac{e^{-i\w_2t}}{2\w_2}\right],
\qquad
W_{zz}(t)\Big|_{\Delta\to0}
=-\,\frac{(1+i\w t)\,e^{-i\w t}}{4\w^3}.
\end{equation}
Both expressions contain only positive frequencies; the first is the
momentum-space spectral function \eqref{eq:Wspectral}, the second its
confluent ($\delta'$) limit.  The commutator parts reproduce the classical
fourth-order Green-function difference
$E_P=[\sin\w_2t/\w_2-\sin\w_1t/\w_1]/\Delta$ with the normalization
$E_P(0)=E_P'(0)=E_P''(0)=0$, $E_P'''(0)=1$, and confluent limit
$(\sin\w t-\w t\cos\w t)/(2\w^3)$.

\begin{theorem}[Bridge theorem]\label{thm:bridge}
The selected PU quantization and the Bateman--Turok spectral quantization
determine the same quasifree two-point functional for all $\Delta\ge0$,
after matching the overall action sign and (in the doubly massless case)
the infrared extension: the selected vacuum's two-point distribution is
\eqref{eq:Wspectral}, whose $\Delta=0$ limit is the Bateman--Turok
$\widetilde W(p)=\theta(p^0)\delta_1'(p^2)$ \cite{BT2026}, realized there
by the Jordan mode doublet $e^{-ipx}$, $(1+2i|p|t)e^{-ipx}$ with null
pairing $[a_1,a_2^{\dg}]=[a_2,a_1^{\dg}]\ne0$.  For $\Delta>0$ this
functional is realized through a positive pseudo-Hermitian completion
($\eta$-inner product); at $\Delta=0$ it is realized through the resonant
Krein/Jordan completion with ghost parity.  We do not assert positivity
with respect to a common involution: ``vacuum functional'' refers to the
normalized quasifree functional determined by the two-point distribution
on the fixed field algebra.
\end{theorem}

\begin{proof}
The first statement is the computation \eqref{eq:Wmode} (machine-verified,
including the commutator normalization which fixes the action-sign
convention: the residual overall sign relative to \cite{BT2026} is exactly
the sign of their perfect-square action).  Quasifree functionals are determined
by their two-point distributions, so equality of \eqref{eq:Wmode} with
the Bateman--Turok mode Wightman function --- both equal to
$(1+i\w t)e^{-i\w t}/(4\w^3)$ in matched conventions --- identifies the
functionals.
\end{proof}

\begin{remark}
(i) The theorem upgrades the resonant-boundary picture: the Krein
formulation is not the residue of a positive theory destroyed by an
ultraviolet particle catastrophe (Section~\ref{sec:decoupling} shows no
such catastrophe exists); it is a different completion of the
\emph{same}, analytically continued vacuum functional, adapted to Jordan
dynamics.  (ii) At $\Delta>0$ the completion ambiguity is invisible in
expectation values but material for which operators are Hermitian; at
$\Delta=0$ positivity is unavailable (no positive invariant metric,
\cite{companion1}) and the Krein completion is the only one.  (iii) What
remains genuinely open at the resonant point is the treatment of the
generalized (Jordan) excitations, where the two completions differ by
construction; ghost parity governs precisely that sector \cite{BT2026}.
\end{remark}

\section{Hadamard structure of the selected vacuum}\label{sec:hadamard}

Let $W^+_m$ denote the Minkowski Klein--Gordon positive-frequency
two-point function and define the divided difference and its confluent
limit
\begin{equation}
W^+_{12}=\frac{W^+_{m_1}-W^+_{m_2}}{m_1^2-m_2^2},
\qquad
W^+_{mm}=-\,\partial_{m^2}W^+_m
\quad(\text{up to the sign convention of \eqref{eq:Wspectral}}),
\end{equation}
which by Theorem~\ref{thm:bridge} are the selected two-point functions.

\begin{theorem}[Fourth-order Hadamard property]\label{thm:hadamard}
Let $P=(\Box+m_1^2)(\Box+m_2^2)$, $d=3$ spatial dimensions.  Then:
\begin{enumerate}
\item $P_xW^+_{12}=P_yW^+_{12}=0$, and
$W^+_{12}(x,y)-W^+_{12}(y,x)=iE_P(x,y)$ with
$E_P=(E_{m_1}-E_{m_2})/\Delta$ the fourth-order commutator function;
\item $\WF(W^+_{12})=\mathcal C^+$, the standard positive-frequency
Hadamard wavefront set; the multiplicity of the characteristic polynomial
changes the order of the light-cone singularity, not its wavefront
directions;
\item the short-distance singularity is logarithmic:
with $W_m^+(x)=\frac{1}{4\pi^2}\frac{mK_1(m\rho)}{\rho}$,
$\rho^2=-x^2+i0x^0$, the mass-independent $\frac{1}{4\pi^2\rho^2}$ term
cancels in the divided difference and
\begin{equation}
W^+_{12}(x)=\frac{1}{8\pi^2}\log\rho+\text{(smooth)}
=\frac{1}{16\pi^2}\log\rho^2+\text{(smooth)},
\end{equation}
the coefficient being universal (mass-independent) and nonzero (both
normalizations are displayed here once; the abstract uses the
$\log\rho^2$ form); the confluent case has the identical coefficient;
\item in the doubly massless case $P=\Box^2$ the infrared extension of
$\theta(p^0)\delta'(p^2)$ requires a scale; different admissible scales
change $W$ by a smooth constant, affecting neither the wavefront set nor
the local Hadamard property --- the bridge of Theorem~\ref{thm:bridge}
presumes matched extension conventions;
\item any two fourth-order spectral Hadamard functionals with the same
commutator differ by a smooth bisolution; the fourth-order parametrix
$H_P=(H_{m_1}-H_{m_2})/\Delta$ (confluently $\partial_{m^2}H_m$) has
$W^+_{12}-H_P$ smooth, so Wick powers $:\!\phi^n\!:_{H_P}$ and covariant
perturbation theory are available.
\end{enumerate}
\end{theorem}

\begin{proof}[Proof sketch]
(1) is verified exactly modewise (both branches and the confluent Jordan
mode), with the Green normalization $E_P'''(0)=1$.  For (2), the
momentum-space form \eqref{eq:Wspectral} is a distribution supported on
the mass shells inside the closed forward cone, giving
$\WF\subseteq\mathcal C^+$ by the standard conic-support argument;
alternatively, divided differences and $\partial_{m^2}$ of the smooth
family $m^2\mapsto W^+_m$ preserve the Hadamard wavefront bound.
Equality follows because the coefficient in (3) is nonzero and the
on-shell $\delta/\delta'$ content is nonvanishing on the cone.  (3) is
the Bessel expansion
$mK_1(m\rho)/\rho=\rho^{-2}+\tfrac{m^2}{2}\log(m\rho/2)+\dots$: the
$\rho^{-2}$ terms cancel and the $\log\rho$ coefficient of the divided
difference is $\big[\tfrac{m_1^2}{2}-\tfrac{m_2^2}{2}\big]/\Delta\cdot
\tfrac{1}{4\pi^2}=\tfrac{1}{8\pi^2}$ (machine-verified, as is the
confluent case).  (4) is the observation that the scale enters only
through $\log(\mu\rho)$, so a change of scale adds
$\log(\mu/\bar\mu)/(8\pi^2)$, a constant.  (5) follows from (1)--(3) by
the usual difference argument: the difference of two such functionals is
a bisolution with vanishing antisymmetric part and empty wavefront set.
\end{proof}

\begin{proposition}[$(\pm)$Hadamard split structure]\label{prop:split}
Let $\phi_1=(\Box+m_2^2)\phi/(m_2^2-m_1^2)$ and
$\phi_2=(\Box+m_1^2)\phi/(m_1^2-m_2^2)$ be the local partial-fraction
fields.  In the selected vacuum, exactly,
\begin{equation}
W_{\phi_1\phi_1}=+\frac{W^+_{m_1}}{\Delta},
\qquad
W_{\phi_2\phi_2}=-\frac{W^+_{m_2}}{\Delta},
\qquad
W_{\phi_1\phi_2}=0 .
\end{equation}
\end{proposition}

The ghost branch carries a Klein--Gordon--Hadamard \emph{distribution
with negative Krein signature} --- not a positive KG state under the
ordinary KG $*$-structure: the Krein signature is intrinsic to the
selected functional, while its local singularity structure is standard on
both branches.  This makes the summary statement precise:

\begin{quote}
\emph{The universal selected vacuum is a fourth-order Hadamard quasifree
functional.  Its apparent infinite particle content relative to the split
two-field vacuum reflects inequivalent complex structures --- the
positive-positive versus positive-negative pairing of the two branches
--- not pathological local ultraviolet singularities.}
\end{quote}

The logarithmic singularity is the expected one for a dimension-zero
field in four spacetime dimensions and is milder than the Klein--Gordon
$1/\sigma$; local Wick composites exist, and derivative observables (in
particular the stress tensor) require the usual local subtractions
relative to $H_P$, with no evident ultraviolet obstruction.  A systematic
treatment along the lines of \cite{Radzikowski,BrunettiFredenhagen,
HollandsWald} for the fourth-order operator is now well-posed.

\section{Discussion}\label{sec:discussion}

\paragraph{What the three geometries settle.}
The separation \eqref{eq:threegeom} resolves several tensions in the
literature at once.  The divergent rapidity $r(k)\sim2\log k$ is real,
but it is metric geometry: it makes the phase-space (observable) metric a
Sobolev pair $H^1\oplus H^{-1}$ \cite{companion2} and the Dyson operator
unboundedly non-unitary in every frame --- yet it creates no quanta,
because it acts inside the vacuum-line stabilizer.  The physical
inequivalence to the naive two-field Fock space is vacuum geometry: an
$O(1)$-per-mode, metric-independent statement
(Theorems~\ref{thm:constancy},~\ref{thm:nogo}).  The resonant
$\Delta\to0$ singularity is dynamical geometry: Jordan blocks and the
death of the positive metric, with the vacuum analytic throughout
(Theorem~\ref{thm:sectors} ff.).  Conflating any two of these produces
false conclusions --- e.g.\ ultraviolet particle catastrophes
($\Lambda^{d+2}$ scalings) or mass-splitting superselection sectors, both
of which are refuted by the exact computations here.

\paragraph{Relation to Bender--Mannheim and Bateman--Turok.}
Theorem~\ref{thm:bridge} shows the two programs select the same state,
from opposite directions: positivity of the inner product modewise
(Bender--Mannheim), spectral positivity of the energy covariantly
(Bateman--Turok).  The difference is the completion, and the difference
matters exactly where the completions are forced apart: at the resonant
point, where positivity fails (the Jordan no-go of \cite{companion1}) and
ghost parity takes over the excitation sector.  A sharp remaining
question is the representation-theoretic status of the generalized Jordan
excitations at $\Delta=0$: the vacuum sector is common, the one-particle
completions are not, and the interacting perfect-square theory of
\cite{BT2026} lives precisely there.

\paragraph{Toward quadratic gravity.}
Applied polarization by polarization, the free-field statements transfer
to the scalar and spin-two PU branches of quadratic gravity in its
conformally flat and linearized regimes: the physical vacuum is
spectral, fourth-order Hadamard, metric-independent, and not normal to
the naive two-field (graviton plus ghost) Fock representation; PT-based unitarity
claims should accordingly be phrased in the dressed representation, and
interacting constructions must fix the representation before the
perturbative expansion.  We leave the interacting analysis, the
stress-tensor subtractions relative to $H_P$, and curved backgrounds to
future work.

\begin{thebibliography}{12}

\bibitem{BM2008PRL}
C.~M.~Bender and P.~D.~Mannheim,
Phys.\ Rev.\ Lett.\ \textbf{100} (2008) 110402.

\bibitem{BM2008PRD}
C.~M.~Bender and P.~D.~Mannheim,
Phys.\ Rev.\ D \textbf{78} (2008) 025022.

\bibitem{BT2026}
S.~Bateman and N.~Turok,
\emph{Escape from Ostrogradsky via hidden ghost parity},
arXiv:2607.00096.

\bibitem{companion1}
\emph{Canonical positive symplectic diagonalization of the Pais--Uhlenbeck
oscillator}, companion paper and verification repository, 2026.

\bibitem{companion2}
\emph{The Pais--Uhlenbeck metric as a minimum-distortion principle, and
the representation problem for the fourth-order field}, companion paper,
2026.

\bibitem{Mostafazadeh2010}
A.~Mostafazadeh,
Int.\ J.\ Geom.\ Meth.\ Mod.\ Phys.\ \textbf{7} (2010) 1191.

\bibitem{vN1939}
J.~von Neumann,
Compos.\ Math.\ \textbf{6} (1939) 1.

\bibitem{Radzikowski}
M.~J.~Radzikowski,
\emph{Micro-local approach to the Hadamard condition in quantum field
theory on curved space-time},
Commun.\ Math.\ Phys.\ \textbf{179} (1996) 529.

\bibitem{BrunettiFredenhagen}
R.~Brunetti and K.~Fredenhagen,
\emph{Microlocal analysis and interacting quantum field theories},
Commun.\ Math.\ Phys.\ \textbf{208} (2000) 623.

\bibitem{HollandsWald}
S.~Hollands and R.~M.~Wald,
\emph{Local Wick polynomials and time ordered products of quantum fields
in curved spacetime},
Commun.\ Math.\ Phys.\ \textbf{223} (2001) 289.

\bibitem{Holdom}
B.~Holdom,
Nucl.\ Phys.\ B \textbf{1000} (2024) 116696;
Nucl.\ Phys.\ B \textbf{1008} (2024) 116472.

\bibitem{ABHT}
M.~Anderson, S.~Bateman, F.~Herzog and N.~Turok,
\emph{Renormalization of a four-derivative theory} and
\emph{Conformally flat limit of quadratic gravity}, 2026, to appear.

\end{thebibliography}

\end{document}
